\chapter{纸质练习四：传递矩阵的最小实现}
\label{paperwork:04}

\section{原题}

\ex{纸质练习四\quad 求传递矩阵的最小实现}{
已知
\[
    G(s)=\begin{bmatrix}
      \dfrac{s+1}{(s+2)(s+3)} & \dfrac{5s+4}{s+5}
    \end{bmatrix}.
\]
求该系统的最小实现。
}

题目来源为《纸质练习答案》第 2 页。

\nt{\textbf{所用知识点}\quad
直接传递矩阵 $D$，公共分母，能观规范 II 形，多输入通道的分子系数排列，最小阶数与实际极点。}

\section{详细解答}

\subsection{读出输入、输出和直接传递项}

$G(s)$ 有一行两列，故 $u=[u_1\ u_2]^\top$ 有两个分量，$y$ 为标量。设实现用 $n$ 个状态，
\[
    \dot x=Ax+Bu,\qquad y=Cx+Du,\qquad
    G(s)=C(sI-A)^{-1}B+D .
\]
$B$ 应有两列，$C$ 应有一行，$D$ 的尺寸固定为 $1\times2$。上式的传递矩阵从\textbf{零初态}求得。因为 $(sI-A)^{-1}$ 随 $s\to\infty$ 趋于零，逐通道取极限可得 $D=\lim_{s\to\infty}G(s)$。

第一通道的分子次数小于分母次数，它的极限为零。第二通道的分子、分母同为一次。先把分子写成分母的整数倍加余数，
\[
    5s+4=5(s+5)-21
    \quad\Longrightarrow\quad
    \frac{5s+4}{s+5}=5-\frac{21}{s+5}.
\]
因此 $D=[0\ 5]$，严格真部分为
\[
    G(s)-D=
    \begin{bmatrix}\dfrac{s+1}{(s+2)(s+3)} &-\dfrac{21}{s+5}\end{bmatrix}.
\]
时间域中，$5u_2$ 直接进入 $y$。余数 $-21$ 决定输入 $u_2$ 的动态部分，稍后进入 $B$ 的第二列。

\subsection{通分并逐次提取系数}

取公共分母 $a(s)=(s+2)(s+3)(s+5)$。先乘前两个因子，
\[
\begin{aligned}
    (s+2)(s+3)&=s^2+5s+6,\\
    a(s)&=(s^2+5s+6)(s+5)=s^3+10s^2+31s+30.
\end{aligned}
\]
第一通道分子补上 $s+5$，第二通道分子补上 $(s+2)(s+3)$，所以
\[
\begin{aligned}
    N_1(s)&=(s+1)(s+5)=s^2+6s+5,\\
    N_2(s)&=-21(s+2)(s+3)=-21s^2-105s-126 .
\end{aligned}
\]
按 $s^0,s^1,s^2$ 的顺序读取每个通道的系数，再把同一幂次的系数排在同一行。
\[
\begin{array}{c|rr}
    \text{幂次}&N_1(s)&N_2(s)\\ \hline
    s^0&5&-126\\
    s^1&6&-105\\
    s^2&1&-21
\end{array}
\qquad\Longrightarrow\qquad
    B=\begin{bmatrix}5&-126\\6&-105\\1&-21\end{bmatrix}.
\]
$B$ 的每一\textbf{列}对应一个输入通道，每一\textbf{行}对应一个分子幂次。把这三行依次记为 $r_0,r_1,r_2$，则 $G-D=(r_0+sr_1+s^2r_2)/a(s)$。

\subsection{构造能观规范 II 形}

对分母 $a(s)=s^3+a_2s^2+a_1s+a_0$，本课程的能观规范 II 形把 $-a_0,-a_1,-a_2$ 从上到下排在 $A$ 的最后一列，次对角线为 $1$，输出取最后一个状态。代入 $a_0=30,a_1=31,a_2=10$，得到
\[
 \begin{gathered}
    A=\begin{bmatrix}0&0&-30\\1&0&-31\\0&1&-10\end{bmatrix},\qquad
    B=\begin{bmatrix}5&-126\\6&-105\\1&-21\end{bmatrix},\\[-1pt]
    C=\begin{bmatrix}0&0&1\end{bmatrix},\qquad
    D=\begin{bmatrix}0&5\end{bmatrix}.
 \end{gathered}
\]
把两列输入的作用逐行写出来，便于核对符号，
\[
\boxed{\begin{aligned}
    \dot x_1&=-30x_3+5u_1-126u_2,\\
    \dot x_2&=x_1-31x_3+6u_1-105u_2,\\
    \dot x_3&=x_2-10x_3+u_1-21u_2,\\
    y&=x_3+5u_2 .
\end{aligned}}
\]

\subsection{逐通道验算传递矩阵}

完整求逆会算出许多用不到的元素。由于输出取 $x_3$，只需求行向量 $L(s)=C(sI-A)^{-1}$。设 $L=[\ell_1\ \ell_2\ \ell_3]$，用 $L(sI-A)=C$ 得
\[
    [\ell_1\ \ell_2\ \ell_3]
    \begin{bmatrix}s&0&30\\-1&s&31\\0&-1&s+10\end{bmatrix}
    =[0\ 0\ 1].
\]
比较三个分量，
\[
    s\ell_1-\ell_2=0,\qquad s\ell_2-\ell_3=0,\qquad
    30\ell_1+31\ell_2+(s+10)\ell_3=1.
\]
前两式给出 $\ell_2=s\ell_1,\ \ell_3=s^2\ell_1$。代入第三式得到
$a(s)\ell_1=1$，所以
\[
    L(s)=\frac{[1\ s\ s^2]}{a(s)}.
\]
现在分别乘 $B$ 的第一列、第二列，并加上相应的直接传递项，
\[
\begin{aligned}
    G_{11}(s)&=L(s)B_{:1}+D_{11}
     =\frac{5+6s+s^2}{a(s)}
     =\frac{(s+1)(s+5)}{(s+2)(s+3)(s+5)}
     =\frac{s+1}{(s+2)(s+3)},\\
    G_{12}(s)&=L(s)B_{:2}+D_{12}
     =\frac{-126-105s-21s^2}{a(s)}+5
     =-\frac{21}{s+5}+5
     =\frac{5s+4}{s+5}.
\end{aligned}
\]
第一行的分子出现 $s+5$ 是通分造成的，它约掉后第一通道仍是二阶。第二行的分子含 $(s+2)(s+3)$，这两个因子同样来自通分。判断整个实现的阶数时，要合并两通道各自的实际极点。

\subsection{证明阶数最小}

先作实际极点检查。第一通道在 $s=-2,-3$ 时分子分别为 $-1,-2$，所以其两个极点均保留。第二通道在 $s=-5$ 时分子为 $5(-5)+4=-21\ne0$，故它还有极点 $-5$。三个位置互异。

任意 $n$ 阶实现都满足
\[
    G(s)-D=C\frac{\operatorname{adj}(sI-A)}{\det(sI-A)}B .
\]
因此每个通道的实际极点必须是 $\det(sI-A)$ 的根。本题至少需要三个互异根，故 $n\ge3$。上述实现恰有三个状态，因而阶数已经最小。

还可用能控和能观矩阵独立验算。设 $b_1,b_2$ 是 $B$ 的两列。只需在能控矩阵 $Q_c=[B\ AB\ A^2B]$ 中找到一个三阶非零子式。取前两列与 $Ab_1$，
\[
\begin{aligned}
    Ab_1&=
    \begin{bmatrix}0&0&-30\\1&0&-31\\0&1&-10\end{bmatrix}
    \begin{bmatrix}5\\6\\1\end{bmatrix}
    =\begin{bmatrix}-30\\-26\\-4\end{bmatrix},\\
    \det[b_1\ b_2\ Ab_1]&=
    \begin{vmatrix}5&-126&-30\\6&-105&-26\\1&-21&-4\end{vmatrix}
    =5(-126)+126(2)+(-30)(-21)=252 .
\end{aligned}
\]
故 $\rank Q_c=3$。逐次右乘 $A$，得到 $CA=[0\ 1\ {-10}]$、$CA^2=[1\ {-10}\ 69]$，所以
\[
    Q_o=\begin{bmatrix}0&0&1\\0&1&-10\\1&-10&69\end{bmatrix},
    \qquad \det Q_o=-1 .
\]
实现既完全能控又完全能观，与极点下界的结论一致。

\subsection{另一种答案：用实际极点直接建模}

分解第一通道的分式，
\[
  \frac{s+1}{(s+2)(s+3)}=-\frac{1}{s+2}+\frac{2}{s+3}.
\]
例如分别乘以 $s+2$、$s+3$，再代 $s=-2$ 和 $s=-3$，便得到系数 $-1$ 和 $2$。第二通道已有 $5-21/(s+5)$。让每个极点对应一个一阶状态，取
\[
\begin{aligned}
   \dot z_1&=-2z_1+u_1,&
   \dot z_2&=-3z_2+u_1,&
   \dot z_3&=-5z_3+u_2,\\
   y&=-z_1+2z_2-21z_3+5u_2 . &&
\end{aligned}
\]
零初态下 $Z_1=U_1/(s+2)$、$Z_2=U_1/(s+3)$、$Z_3=U_2/(s+5)$，代入输出即还原 $G(s)$。这个对角实现也是三阶，所以同样是最小实现。检验维数与传递矩阵即可比较不同的答案。

\section{同类题怎么做}

\subsection{解题顺序与易错检查}

\begin{enumerate}
    \item 先数 $G(s)$ 的行与列，再对每个元素长除并约分。$D$ 的尺寸与 $G$ 相同。
    \item 从约分后的严格真部分取首一公共分母。一个输出时能观规范 II 形的初始维数等于公分母次数，一个输入时可优先用能控规范 II 形。
    \item 把公共分母系数放入 $A$，按升幂把每个分子放入 $B$ 的相应列。核对 $y=Cx+Du$ 中的直接项，最后回算每个通道。
    \item 计算能控、能观秩，或找一个能达到的状态数下界。具有多个输出的传递矩阵还可能在同一个极点需要多个独立的模态方向。
\end{enumerate}

\begin{note}
    对一般 MIMO 系统，互异极点数只是下界。例如
    $G(s)=I_2/(s+1)$ 的两个通道都只有 $-1$，但该极点的留数矩阵是 $I_2$，秩为 $2$。若只有一个状态，该留数只能写成一列乘一行，秩至多为 $1$。取 $A=-I_2,\ B=C=I_2,\ D=0$ 可用两状态实现。
\end{note}

\subsection{小变式：约分后两路输入共用一个状态}

\ex{变式\quad 两路输入共享同一个极点}{
求
\[
    \widehat G(s)=\begin{bmatrix}
        \dfrac{s+2}{(s+2)(s+3)} & \dfrac{5s+4}{s+3}
    \end{bmatrix}
\]
的最小实现。
}

\sol
第一通道先消去 $s+2$。第二通道有 $5s+4=5(s+3)-11$，所以
\[
   \widehat G(s)=[0\ 5]+\frac{[1\ {-11}]}{s+3}.
\]
共用一个状态即可，取
\[
   \dot z=-3z+u_1-11u_2,\qquad y=z+5u_2 .
\]
验算 $Z=(U_1-11U_2)/(s+3)$，得到所给 $\widehat G(s)$。其中实际极点 $-3$ 要求至少一个状态，所以该实现最小。

若沿用未约分的分母 $(s+2)(s+3)$，会得到一个二阶能观规范 II 形
\[
   \widehat A=\begin{bmatrix}0&-6\\1&-5\end{bmatrix},\quad
   \widehat B=\begin{bmatrix}2&-22\\1&-11\end{bmatrix},\quad
   \widehat C=[0\ 1],\quad\widehat D=[0\ 5].
\]
$\widehat B$ 的两列成比例，且 $\widehat A[2\ 1]^\top=[-6\ {-3}]^\top=-3[2\ 1]^\top$，故其能控秩只有 $1$。少掉的一维正是约分消去的 $-2$ 模态。本题先约分可以直接得到最低阶。
